%%% Local Variables:
%%% mode: LaTeX
%%% TeX-engine: luatex
%%% TeX-master: t
%%% TeX-backend: "biber"
%%% End:
\documentclass[book,spanish,alphabetic]{apuntes-scr}

\subject{Manual de Referencia Oficial}

\usepackage{tabularx}
\usepackage{booktabs}

\title{Manual de Referencia Oficial e Integral de la Clase \texttt{apuntes-scr.cls}}
\subtitle{Guía Definitiva de Arquitectura Tipográfica, Estructura Bourbaki / IHÉS y Redacción Científica}
\author{Campos Mendoza, D. Daniel}
\date{%
    \vspace{2em}
    \textbf{Ecosistema GNU Emacs 30 / Sway / Wayland} \\
    \vspace{0.5em}\normalsize
    \textbf{Fecha:} \today \quad \textbf{Hora:} \DTMcurrenttime
}

\begin{document}

\frontmatter
\maketitle
\tableofcontents

\mainmatter

\chapter{Introducción y Filosofía de Diseño}

La clase \texttt{apuntes-scr.cls} es un sistema tipográfico de nivel editorial diseñado para la redacción de textos matemáticos de alta exigencia, monografías científicas, apuntes de posgrado y libros de investigación. Construida sobre la base sólida de \textbf{KOMA-Script} (\texttt{scrbook} y \texttt{scrartcl}), está optimizada específicamente para el motor \textbf{Lua\LaTeX} con soporte nativo de \textbf{HarfBuzz}, \textbf{OpenType} y fuentes matemáticas \textbf{Unicode-Math}.

\section{Pilares Editoriales de la Clase}
\begin{enumerate}[label=\bfseries\arabic*., leftmargin=10mm]
    \item \textbf{Estética Bourbaki / IHÉS:} Inspirada directamente en los \textit{Éléments de mathématique} de Nicolas Bourbaki y las publicaciones del \textit{Institut des Hautes Études Scientifiques} (IHÉS). Utiliza proporciones geométricas áureas, filetes sobrios y márgenes exteriores generosos para notas y aclaraciones marginales.
    \item \textbf{Tipografía Monumental STIX Two:} Combina \textbf{STIX Two Text} y \textbf{STIX Two Math} con conjuntos estilísticos avanzados (\textit{Stylistic Sets} 1 y 8), ofreciendo máxima legibilidad en expresiones algebraicas complejas, símbolos de teoría de esquemas y categorías.
    \item \textbf{Sistema Completo de Entornos Vectoriales (\texttt{tcolorbox}):} Todos los teoremas, definiciones, lemas y problemas cuentan con recuadros vectoriales nítidos con numeración unificada por capítulo y compatibilidad total con \texttt{cleveref}.
    \item \textbf{Modularidad Total (Segundo Cerebro):} Gracias a su integración con \texttt{docmute} e \texttt{import}, cualquier subdocumento escrito con \texttt{apuntes-scr} puede compilarse de forma autónoma o integrarse directamente dentro del libro maestro unificado sin colisiones de preámbulo.
\end{enumerate}

\chapter{Opciones de Clase y Plantillas de Inicio}

\section{Opciones Globales Disponibles}
Al invocar la clase:
\begin{lstlisting}[language=TeX, basicstyle=\ttfamily\footnotesize, frame=single, backgroundcolor=\color{gray!5}]
\documentclass[opciones]{apuntes-scr}
\end{lstlisting}
se pueden especificar las siguientes opciones:

\begin{center}
\small
\begin{tabularx}{\linewidth}{p{2.5cm} p{3.2cm} X}
\toprule
\textbf{Categoría} & \textbf{Opción} & \textbf{Comportamiento y Efecto} \\
\midrule
\textbf{Idioma} & \texttt{spanish} (defecto) & Español con reglas tipográficas canónicas y signos abiertos. \\
                & \texttt{english} & Inglés internacional con capitalización en entornos. \\
                & \texttt{french} & Francés según normas de la \textit{Imprimerie Nationale}. \\
                & \texttt{german} & Alemán con ortografía moderna (\texttt{ngerman}). \\
\midrule
\textbf{Estructura} & \texttt{book} (defecto) & Basado en \texttt{scrbook}, con capítulos (\texttt{\string\chapter}), partes y front/backmatter. \\
                    & \texttt{article} & Basado en \texttt{scrartcl}, para artículos, exámenes o notas cortas sin capítulos. \\
\midrule
\textbf{Bibliografía} & \texttt{alphabetic} (def) & Citas estilo Bourbaki (ej. [Bou06], [Ati69]). \\
                      & \texttt{numeric} & Citas numéricas entre corchetes (ej. [1], [2]). \\
                      & \texttt{ieee}, \texttt{apa} & Estilos formales de ingeniería y ciencias sociales. \\
\midrule
\textbf{Rendimiento} & \texttt{fast} & Modo borrador rápido que desactiva compilaciones pesadas de código y optimiza el procesamiento de figuras. \\
\bottomrule
\end{tabularx}
\end{center}

\section{Plantilla Mínima para un Libro de Investigación}
\begin{lstlisting}[language=TeX, basicstyle=\ttfamily\footnotesize, frame=single, backgroundcolor=\color{gray!5}]
\documentclass[book,spanish,alphabetic]{apuntes-scr}

\subject{Manual de Referencia Oficial}

\title{Geometría Algebraica y Esquemas}
\subtitle{Notas de Posgrado}
\author{Campos Mendoza, D. Daniel}
\date{%
    \vspace{2em}
    \textbf{Ecosistema GNU Emacs 30 / Sway / Wayland} \\
    \vspace{0.5em}\normalsize
    \textbf{Fecha:} \today \quad \textbf{Hora:} \DTMcurrenttime
}

\begin{document}

\frontmatter
\maketitle
\tableofcontents

\mainmatter

\chapter{Espectro de un Anillo y Haces Estructurales}
\section{Topología de Zariski}

\backmatter
\printindex

\end{document}
\end{lstlisting}

\section{Plantilla para Artículo / Evaluación Universitaria (\texttt{article})}
\begin{lstlisting}[language=TeX, basicstyle=\ttfamily\footnotesize, frame=single, backgroundcolor=\color{gray!5}]
\documentclass[spanish,article,numeric]{apuntes-scr}

\title{Práctica Calificada 1: Teoría de Probabilidad}
\author{Campos Mendoza, D. Daniel}
\date{%
    \vspace{2em}
    \textbf{Ecosistema GNU Emacs 30 / Sway / Wayland} \\
    \vspace{0.5em}\normalsize
    \textbf{Fecha:} \today \quad \textbf{Hora:} \DTMcurrenttime
}

\begin{document}
\maketitle

\section{Problema 1: Martingalas en Tiempo Discreto}

\end{document}
\end{lstlisting}

\chapter{Sistema de Teoremas y Cajas Matemáticas}

La clase implementa un catálogo exhaustivo de entornos con numeración unificada por capítulo y colores armónicos:

\section{Entornos Teóricos Principales}

\begin{definition}[Espacio Anillado]\label{def:espacio_anillado}
    \index{Espacio anillado} Un \deftech{espacio anillado} es un par $(X, \symcal{O}_X)$ formado por un espacio topológico $X$ y un haz de anillos $\symcal{O}_X$ sobre $X$. Si para cada $x \in X$ el tallo $\symcal{O}_{X,x}$ es un anillo local, decimos que $(X, \symcal{O}_X)$ es un \deftech{espacio localmente anillado}.
\end{definition}

\begin{theorem}[Teorema de Nakayama]\label{thm:nakayama}
    Sea $A$ un anillo conmutativo unitario, $\symfrak{m} \subseteq A$ un ideal contenido en el radical de Jacobson $\Rad(A)$, y $M$ un $A$-módulo finitamente generado. Si $\symfrak{m} M = M$, entonces $M = 0$.
\end{theorem}
\begin{proof}
    Supongamos por reducción al absurdo que $M \neq 0$. Sea $\{u_1, \dots, u_n\}$ un sistema minimal de generadores de $M$ con $n \ge 1$. Dado que $u_n \in M = \symfrak{m} M$, existen elementos $a_1, \dots, a_n \in \symfrak{m}$ tales que:
    \[
        u_n = \sum_{i=1}^n a_i u_i \implies (1 - a_n) u_n = \sum_{i=1}^{n-1} a_i u_i.
    \]
    Como $a_n \in \symfrak{m} \subseteq \Rad(A)$, el elemento $1 - a_n$ es una unidad en $A$. Multiplicando por $(1 - a_n)^{-1}$ se obtiene que $u_n$ pertenece al submódulo generado por $\{u_1, \dots, u_{n-1}\}$, contradiciendo la minimalidad del sistema generador.
\end{proof}

\begin{lemma}[Lema de Urysohn]\label{lem:urysohn}
    Sea $X$ un espacio topológico normal y sean $A, B \subseteq X$ dos subconjuntos cerrados disjuntos. Entonces existe una función continua $f \colon X \to [0, 1]$ tal que $f(A) = \{0\}$ y $f(B) = \{1\}$.
\end{lemma}

\begin{proposition}[Caracterización de Módulos Proyectivos]\label{prop:mod_proyectivos}
    Para un $A$-módulo $P$, las siguientes afirmaciones son equivalentes:
    \begin{tfae}
        \item $P$ es proyectivo.
        \item El funtor $\Hom_A(P, -)$ es exacto.
        \item Toda sucesión exacta corta $0 \to M' \to M \to P \to 0$ se escinde.
        \item $P$ es un sumando directo de un módulo libre.
    \end{tfae}
\end{proposition}

\begin{corollary}[Corolario de Libertad Local]\label{cor:libertad_local}
    Todo módulo finitamente generado y proyectivo sobre un anillo local es libre.
\end{corollary}

\section{Entornos Descriptivos y Ejemplos}

\begin{example}[El Haz Estructural de una Variedad Afín]\label{ex:haz_afin}
    Sea $X = \Spec(A)$ un esquema afín. Para cada abierto básico $D(f)$ con $f \in A$, las secciones globales del haz estructural $\symcal{O}_X(D(f))$ coinciden exactamente con la localización $A_f$.
\end{example}

\begin{counterexample}[Un Anillo No Noetheriano]\label{cex:no_noetheriano}
    El anillo de polinomios en una cantidad infinita numerable de indeterminadas $A = k[X_1, X_2, \dots]$ no es noetheriano, pues la cadena ascendente de ideales:
    \[
        (X_1) \subsetneq (X_1, X_2) \subsetneq (X_1, X_2, X_3) \subsetneq \cdots
    \]
    no se estabiliza jamás.
\end{counterexample}

\begin{remark}[Observación Metodológica]\label{rem:metodo}
    Nótese que en teoría de haces es crucial distinguir entre el tallo $\symcal{F}_x$ en un punto y la fibra sobre dicho punto.
\end{remark}

\begin{notation}\label{not:notacion_estandar}
    Denotaremos por $k[X_1, \dots, X_n]$ el álgebra de polinomios sobre un cuerpo base $k$ y por $\symbb{A}_k^n = \Spec(k[X_1, \dots, X_n])$ el espacio afín de dimensión $n$.
\end{notation}

\section{Afirmaciones y Demostraciones Locales}

\begin{claim}[Inyectividad de la Restricción]\label{clm:inyectividad}
    La aplicación de restricción $\rho_{U, V} \colon \symcal{F}(U) \to \symcal{F}(V)$ es inyectiva para todo abierto conexo $V \subseteq U$.
\end{claim}
\begin{claimproof}
    Sea $s \in \symcal{F}(U)$ tal que $s|_V = 0$. Dado que $\symcal{F}$ es un haz separable y $X$ es conexo, la sección se anula idénticamente en todo $U$.
\end{claimproof}

\section{Problemas, Ejercicios y Soluciones}

\begin{exercise}[Cálculo del Grupo Fundamental]\label{exer:grupo_fundamental}
    Calcular el grupo fundamental del toro bidimensional $\symbb{T}^2 = S^1 \times S^1$ y describir sus generadores geométricos.
\end{exercise}
\begin{solution}
    Dado que el toro es el producto cartesiano de dos circunferencias conexas por caminos y localmente contractibles:
    \[
        \pi_1(\symbb{T}^2) \cong \pi_1(S^1) \times \pi_1(S^1) \cong \Z \times \Z.
    \]
    Los generadores canónicos corresponden a los lazos meridianos y paralelos que rodean cada factor del producto.
\end{solution}

\chapter{Estructuras de Razonamiento Lógico e Inducción}

\section{Pasos Formales de Demostración}
Para estructurar pruebas rigurosas paso a paso:
\begin{itemize}[leftmargin=8mm]
    \item \texttt{\string\directstep} genera: \directstep Demostración de la implicación directa ($\Rightarrow$).
    \item \texttt{\string\reversestep} genera: \reversestep Demostración del recíproco ($\Leftarrow$).
    \item \texttt{\string\containedstep} genera: \containedstep Demostración de la contención directa ($\subset$).
    \item \texttt{\string\inversecontainedstep} genera: \inversecontainedstep Demostración de la contención inversa ($\supset$).
\end{itemize}

\section{Inducción Matemática Formal}
La clase provee directivas semánticas para pruebas por inducción:

\casobase{n=1} Para $n=1$, la igualdad se verifica inmediatamente:
\[
    \sum_{i=1}^1 i = 1 = \frac{1(1+1)}{2}.
\]

\hipotesisind Supongamos cierta la fórmula para $n = k$, es decir, $\sum_{i=1}^k i = \frac{k(k+1)}{2}$.

\pasoinductivo Para $n = k+1$, aplicando la hipótesis de inducción:
\[
    \sum_{i=1}^{k+1} i = \left(\sum_{i=1}^k i\right) + (k+1) = \frac{k(k+1)}{2} + (k+1) = \frac{(k+1)(k+2)}{2}.
\]
Esto concluye el paso inductivo.

\section{Párrafos Numerados Estilo EGA (\textit{Éléments de Géométrie Algébrique})}
Para organizar tratados monumentales mediante artículos numerados de forma compacta:

\numpar[Morfismos Separados] Un morfismo de esquemas $f \colon X \to Y$ se dice \deftech{separado} si el morfismo diagonal $\Delta_{X/Y} \colon X \to X \times_Y X$ es una inmersión cerrada.

\numpar[Propiedades Topológicas] Si $Y$ es afín y $f$ es separado, entonces la intersección de dos abiertos afines en $X$ es nuevamente un abierto afín.

\chapter{Cajas de Información, Consejos y Resumen}

\section{Cajas Especiales de Alerta y Estudio}

\begin{notabox}[Consejo de Estudio]
    Para asimilar los conceptos de álgebra homológica, es recomendable verificar cada propiedad universal dibujando explícitamente el diagrama conmutativo correspondiente antes de iniciar el cálculo algebraico.
\end{notabox}

\begin{warningbox}[Advertencia Importante]
    La aplicación de localización no preserva en general ideales maximales; la imagen de un ideal maximal puede degenerar en el anillo total en el cuerpo de fracciones.
\end{warningbox}

\begin{controlbox}[Pregunta de Autoevaluación]
    ¿Por qué el lema de Nakayama exige que el módulo $M$ sea finitamente generado? ¿Qué contraejemplo surge considerando $M = \Q$ sobre $A = \Z_{(p)}$?
\end{controlbox}

\section{Marcadores de Estado en la Investigación}
\begin{itemize}[leftmargin=8mm]
    \item \texttt{\string\deftech\{Variedad afín\}}: Resalta el concepto e indexa automáticamente en el índice analítico.
    \item \texttt{\string\TODO\{Completar demostración del lema 3\}}: \TODO{Completar demostración del lema 3}.
    \item \texttt{\string\DEBUG\{Verificar compatibilidad de signos\}}: \DEBUG{Verificar compatibilidad de signos}.
    \item \texttt{\string\NOTE\{Revisar versión de Hartshorne Cap. II\}}: \NOTE{Revisar versión de Hartshorne Cap. II}.
\end{itemize}

\chapter{Diccionario Canónico de Operadores Matemáticos}

La clase estandariza más de 60 operadores matemáticos para evitar el uso de comandos planos como \texttt{\string\text\{Hom\}}:

\section{Conjuntos Numéricos y Cuerpos}
\begin{center}
\small
\begin{tabularx}{\linewidth}{l l X}
\toprule
\textbf{Macro} & \textbf{Salida} & \textbf{Significado Matemático} \\
\midrule
\texttt{\string\C} & $\C$ & Números Complejos ($\symbb{C}$) \\
\texttt{\string\R} & $\R$ & Números Reales ($\symbb{R}$) \\
\texttt{\string\Q} & $\Q$ & Números Racionales ($\symbb{Q}$) \\
\texttt{\string\Z} & $\Z$ & Números Enteros ($\symbb{Z}$) \\
\texttt{\string\N} & $\N$ & Números Naturales ($\symbb{N}$) \\
\texttt{\string\K} & $\K$ & Cuerpo base arbitrario ($\symbb{K}$) \\
\texttt{\string\F} & $\F$ & Cuerpo finito o general ($\symbb{F}$) \\
\texttt{\string\symbb\{T\}} & $\symbb{T}$ & Toro algebraico / complejo ($\symbb{T}$) \\
\texttt{\string\symbb\{P\}} & $\symbb{P}$ & Espacio proyectivo ($\symbb{P}$) \\
\texttt{\string\symbb\{A\}} & $\symbb{A}$ & Espacio afín ($\symbb{A}$) \\
\bottomrule
\end{tabularx}
\end{center}

\section{Álgebra Abstracta, Categorías y Álgebras de Lie}
\begin{center}
\small
\begin{tabularx}{\linewidth}{l l X}
\toprule
\textbf{Macro} & \textbf{Salida} & \textbf{Significado Matemático} \\
\midrule
\texttt{\string\Hom(A, B)} & $\Hom(A, B)$ & Conjunto / Módulo de homomorfismos \\
\texttt{\string\End(V)} & $\End(V)$ & Álgebra de endomorfismos \\
\texttt{\string\Aut(G)} & $\Aut(G)$ & Grupo de automorfismos \\
\texttt{\string\Ext\string^{i}(M, N)} & $\Ext^i(M, N)$ & Funtor Ext de extensiones homológicas \\
\texttt{\string\Tor\string_{i}(M, N)} & $\Tor_i(M, N)$ & Funtor Tor de torsión \\
\texttt{\string\Ker(f)} & $\Ker(f)$ & Núcleo de una aplicación lineal \\
\texttt{\string\Coker(f)} & $\Coker(f)$ & Conúcleo \\
\texttt{\string\im(f)} & $\im(f)$ & Imagen \\
\texttt{\string\GL(V)} & $\GL(V)$ & Grupo lineal general \\
\texttt{\string\SL(V)} & $\SL(V)$ & Grupo lineal especial \\
\texttt{\string\mathrm\{SO\}(n)} & $\mathrm{SO}(n)$ & Grupo ortogonal especial \\
\texttt{\string\mathrm\{Sp\}(2n)} & $\mathrm{Sp}(2n)$ & Grupo simpléctico \\
\texttt{\string\glie\{V\}} & $\glie{V}$ & Álgebra de Lie $\symfrak{gl}(V)$ \\
\texttt{\string\Tr(M)} & $\Tr(M)$ & Traza de una matriz u operador \\
\texttt{\string\Span(S)} & $\Span(S)$ & Espacio generado \\
\bottomrule
\end{tabularx}
\end{center}

\section{Geometría Algebraica y Álgebra Conmutativa}
\begin{center}
\small
\begin{tabularx}{\linewidth}{l l X}
\toprule
\textbf{Macro} & \textbf{Salida} & \textbf{Significado Matemático} \\
\midrule
\texttt{\string\Spec(A)} & $\Spec(A)$ & Espectro primo de un anillo \\
\texttt{\string\Specmax(A)} & $\Specmax(A)$ & Espectro maximal \\
\texttt{\string\Proj(S)} & $\Proj(S)$ & Esquema proyectivo de un anillo graduado \\
\texttt{\string\Supp(M)} & $\Supp(M)$ & Soporte de un módulo o haz \\
\texttt{\string\Ann(M)} & $\Ann(M)$ & Anulador de un módulo \\
\texttt{\string\Rad(I)} & $\Rad(I)$ & Radical de un ideal \\
\texttt{\string\nil(A)} & $\nil(A)$ & Nilradical de un anillo \\
\texttt{\string\Pic(X)} & $\Pic(X)$ & Grupo de Picard \\
\texttt{\string\Frac(A)} & $\Frac(A)$ & Cuerpo de fracciones \\
\bottomrule
\end{tabularx}
\end{center}

\section{Análisis, Cálculo y Formas Diferenciales}
\begin{itemize}[leftmargin=8mm]
    \item \textbf{Diferenciales:} \texttt{\string\diff x \string\wedge \string\diff y} $\to \diff x \wedge \diff y$.
    \item \textbf{Parte real e imaginaria:} \texttt{\string\re\{z\} + i \string\im\{z\}} $\to \re{z} + i \im{z}$.
    \item \textbf{Normas y Productos:} \texttt{\string\| v \string\|} $\to \|v\|$, \texttt{\string\langle u, v \string\rangle} $\to \langle u, v \rangle$.
    \item \textbf{Operadores Vectoriales:} \texttt{\string\divop(\string\vec\{F\})} $\to \divop(\vec{F})$.
\end{itemize}

\chapter{Estructuras de Algoritmos y Listas Bourbaki}

\section{Listas de Equivalencia Lógica (\texttt{tfae})}
Para enunciar teoremas con condiciones equivalentes:
\begin{lstlisting}[language=TeX, basicstyle=\ttfamily\footnotesize, frame=single, backgroundcolor=\color{gray!5}]
\begin{tfae}
    \item $A$ es un dominio euclídeo.
    \item $A$ es un anillo de ideales principales (DIP).
    \item $A$ es un dominio de factorización única (DFU).
\end{tfae}
\end{lstlisting}

\section{Entorno Editorial de Algoritmos}

\begin{algorithm}[Algoritmo de Euclides Extendido]\label{alg:euclides}
    \alginput{Dos enteros positivos $a, b \in \Z$ con $a \ge b$.}
    \algoutput{El máximo común divisor $d = \gcd(a, b)$ y los coeficientes de Bézout $x, y \in \Z$ tales que $a x + b y = d$.}
    \alginit{$(r_0, r_1) \leftarrow (a, b)$, $(s_0, s_1) \leftarrow (1, 0)$, $(t_0, t_1) \leftarrow (0, 1)$.}
    \algstop{Cuando el residuo sea nulo: $r_{k+1} = 0$.}

    \begin{algsteps}
        \item \textbf{División entera:} Calcular el cociente $q_k = \lfloor r_{k-1} / r_k \rfloor$ y el residuo $r_{k+1} = r_{k-1} - q_k r_k$.
        \item \textbf{Actualización de coeficientes:}
        \begin{align*}
            s_{k+1} &\leftarrow s_{k-1} - q_k s_k, \\
            t_{k+1} &\leftarrow t_{k-1} - q_k t_k.
        \end{align*}
        \item Si $r_{k+1} \neq 0$, incrementar $k \leftarrow k+1$ y volver al Paso 1.
        \item \textbf{Retorno final:} Retornar $d = r_k$, $x = s_k$, $y = t_k$.
    \end{algsteps}
\end{algorithm}

\chapter{Gráficos Vectoriales, Diagramas y Código}

\section{Diagramas Conmutativos con TikZ-CD}
La clase incluye soporte preconfigurado para álgebra homológica:

\begin{center}
\begin{tikzcd}[row sep=large, column sep=large]
0 \arrow[r] & \symfrak{a} \arrow[r, "\iota"] \arrow[d, "\phi|_{\symfrak{a}}"] & A \arrow[r, "\pi"] \arrow[d, "\phi"] & A/\symfrak{a} \arrow[r] \arrow[d, dashed, "\bar{\phi}"] & 0 \\
0 \arrow[r] & \symfrak{b} \arrow[r, "\iota'"] & B \arrow[r, "\pi'"] & B/\symfrak{b} \arrow[r] & 0
\end{tikzcd}
\end{center}

\section{Inserción de Figuras Inkscape (\texttt{pdf\_tex})}
Para garantizar que las figuras conserven exactamente la misma tipografía matemática que el texto del documento:

\begin{lstlisting}[language=TeX, basicstyle=\ttfamily\footnotesize, frame=single, backgroundcolor=\color{gray!5}]
\begin{figure}[htpb]
    \centering
    \def\svgwidth{0.75\linewidth}
    \subimport*{figuras/}{esfera-riemann.pdf_tex}
    \caption{Proyección estereográfica de la Esfera de Riemann $\symbb{P}^1(\symbb{C})$.}
    \label{fig:esfera_riemann}
\end{figure}
\end{lstlisting}

\chapter{Referencias Cruzadas, Índices y Segundo Cerebro}

\section{Referencias Cruzadas Inteligentes con \texttt{cleveref}}
La clase tiene configuradas todas las declinaciones en español. No es necesario escribir la palabra \textit{Teorema}, \textit{Lema} o \textit{Ecuación}:
\begin{itemize}[leftmargin=8mm]
    \item \texttt{\string\cref\{thm:nakayama\}} genera: \cref{thm:nakayama}.
    \item \texttt{\string\Cref\{def:espacio\_anillado\}} genera: \Cref{def:espacio_anillado}.
    \item \texttt{\string\cref\{alg:euclides\}} genera: \cref{alg:euclides}.
\end{itemize}

\section{Arquitectura Modular del Segundo Cerebro}
Para integrar apuntes independientes en un libro unificado:
\begin{enumerate}[label=\bfseries\arabic*., leftmargin=10mm]
    \item Cada tema matemático (ej. Álgebra Conmutativa, Superficies de Riemann) mantiene su propio archivo \texttt{.tex} autónomo con \texttt{\string\documentclass\{apuntes-scr\}}.
    \item En el archivo maestro \texttt{Segundo-Cerebro/main.tex}, se carga la clase y se incluyen los archivos con:
    \begin{lstlisting}[language=TeX, basicstyle=\ttfamily\footnotesize, frame=single, backgroundcolor=\color{gray!5}]
\import{10-Matematicas/Apuntes-Algebra-Conmutativa/}{Apuntes-Algebra-Conmutativa.tex}
    \end{lstlisting}
    \item El paquete \texttt{docmute} gestiona de forma transparente los preámbulos para que el libro compile como una sola obra armónica de cientos de páginas con índice general e índice analítico unificados.
\end{enumerate}

\backmatter
\printindex

\end{document}
